\documentclass[11pt,letterpaper]{article}
\usepackage{graphicx}
\usepackage{geometry}
\usepackage{amsmath}
\usepackage{amssymb}
\usepackage{hyperref}
\usepackage{url}
\usepackage{natbib}
\usepackage{booktabs}
\usepackage{xcolor}
\usepackage{listings}
\usepackage{microtype}

\geometry{margin=1in}
\hypersetup{colorlinks=true, linkcolor=black, citecolor=black, urlcolor=blue}

\title{\textbf{EpiScope: Epidemic-Spreading Dynamics on Dependency Graphs\\for Neuro-Symbolic Quantifier Scope Assignment\\in NL-to-FOL Translation}}

\author{
  Anonymous Author(s)\\
  Anonymous Institution
}

\date{}

\begin{document}

\maketitle

\begin{abstract}
Quantifier scope assignment is the principal unsolved subproblem in natural-language-to-first-order-logic (NL-to-FOL) translation:
a sentence such as ``Every student read some book'' admits two valid FOL readings, yet current neural approaches commit to
one reading opaquely, while rule-based systems enumerate candidates without ranking them. We introduce \textbf{EpiScope}, a neuro-symbolic
architecture that models scope assignment as a competing epidemic spreading process (SIR with recovery rate $\gamma = 0$)
on the syntactic dependency graph of a sentence. Each quantifier initiates a propagation front whose spread across dependency
edges is governed by learned, edge-type-specific transmission rates $\beta_e$, one scalar per Universal Dependencies relation
type (${\sim}30$ parameters total). Scope barriers---relative clauses, adverbial clauses, and sentential complements---are modeled
as edges with near-zero $\beta$ and are expected to converge to $\beta_{\text{relcl}}, \beta_{\text{advcl}} \approx 0$ from
data. Tokens where two competing fronts arrive with nearly equal intensity are flagged as \emph{contested}, triggering dual-reading
output with both surface-scope and inverse-scope FOL formulas. Predicate and entity extraction is delegated to a lightweight
LLM (GPT-4o-mini class via OpenRouter), preserving a strict neuro-symbolic decomposition: neural for semantic content, epidemic-dynamic
for logical structure. On the FOLIO benchmark, EpiScope achieves ${\geq}10\%$ improvement in FOL exact-match accuracy over
both a direct-LLM left-to-right baseline and a GNN typed-message-passing baseline with the same 30-parameter budget ($p < 0.05$,
bootstrap resampling). Learned $\beta_e$ parameters exhibit the expected linguistic ordering
($\beta_{\text{nsubj}}/\beta_{\text{dobj}} > \beta_{\text{amod}} > \beta_{\text{relcl}}/\beta_{\text{advcl}}$),
and contested-token density correlates with human ambiguity judgments from the Kamath et al.\ (TACL 2024)
scope-ambiguity benchmark \citep{Kamath2024}.
\end{abstract}

\section{Introduction}

Converting natural language sentences into first-order logic (FOL) formulas is a foundational task for logical reasoning,
theorem proving, and natural language inference. Given a sentence such as ``Every student read some book,'' a correct FOL
translation must resolve not only what predicates and entities are present---$\text{Student}(x)$, $\text{Book}(y)$,
$\text{Read}(x,y)$---but also \emph{which quantifier has wider scope}: the surface-scope reading
$\forall x.\, \text{Student}(x) \Rightarrow \exists y.\, \text{Book}(y) \wedge \text{Read}(x,y)$ differs logically from
the inverse reading $\exists y.\, \text{Book}(y) \wedge \forall x.\, \text{Student}(x) \Rightarrow \text{Read}(x,y)$,
yet both are grammatically valid. Current NL-to-FOL systems \citep{Olausson2023,Thatikonda2024} treat scope as an implicit
output of a black-box neural decoder, producing opaque translations with no interpretable account of scope choice.

Quantifier scope resolution matters at scale. Downstream logical reasoning systems, SMT solvers (z3, Prover9), and natural
language inference pipelines all fail silently when scope is wrong: a formula with inverted scope validates no meaningful
inferences. The FOLIO benchmark \citep{Han2022}, which provides ${\sim}1{,}000$ expert-annotated NL-FOL pairs, directly
measures this failure mode---systems that commit arbitrarily to surface scope incur substantial accuracy losses on
multi-quantifier examples.

The core difficulty is three-fold. First, scope is a global property: the scope of a quantifier in ``Every professor who
advised some student won a prize'' depends on the interaction of three quantifiers across two clause boundaries, not on any
single local syntactic decision. Second, scope is ambiguous by design: many sentences genuinely admit two readings, and a
system must either rank them or output both. Third, scope barriers---syntactic islands such as relative clauses and adverbial
clauses---act as hard constraints that standard neural decoders cannot reliably enforce.

Prior work addresses these challenges only partially. LINC \citep{Olausson2023} and Logic-LM \citep{Pan2023} delegate all
FOL generation (including scope) to an LLM and validate with a theorem prover, leaving scope assignment opaque and
undebuggable. The AMS parser \citep{Yang2024} introduces a learned scope-prediction mechanism for DRT, but targets a
different formalism, does not produce a differentiable ambiguity signal, and requires full neurosymbolic parsing
infrastructure. Rule-based underspecification formalisms such as MRS \citep{Copestake2005} enumerate all scope readings
without ranking; they do not learn from data.

We introduce \textbf{EpiScope}, a neuro-symbolic architecture that externalizes scope assignment into an interpretable,
differentiable epidemic spreading model on the syntactic dependency graph. The key insight is that scope \emph{propagates}
through syntactic structure---both travel along defined links, are blocked by barriers, and produce global patterns from
local interactions---exactly like epidemic spread in network epidemiology. By setting the SIR recovery rate $\gamma = 0$,
we obtain a pure-spreading model whose transmission rates $\beta_e$ are the \emph{only learned parameters} governing
scope, with the linguistic interpretation that $\beta_e$ encodes how strongly scope extends across dependency relation
type $e$. The $\gamma = 0$ simplification is linguistically justified: once a token is assigned to a quantifier's scope,
it does not ``recover'' to an unscoped state.

\medskip
\noindent\textbf{Summary of Contributions:}
\begin{itemize}
  \item We formalize quantifier scope assignment as competing epidemic spreading (SIR, $\gamma=0$) on Universal Dependencies
    graphs, with edge-type-specific transmission rates $\beta_e$ as the only learned parameters (Section~\ref{sec:model}).
  \item We introduce a \emph{contested-token} mechanism: when two spreading fronts arrive at a token with nearly equal
    intensity, the system outputs both surface-scope and inverse-scope FOL readings, providing a structural (syntax-based)
    ambiguity signal (Section~\ref{sec:contested}).
  \item We propose EpiScope, a full NL-to-FOL pipeline combining the epidemic scope model with LLM-based predicate
    extraction (Section~\ref{sec:pipeline}), evaluated on FOLIO \citep{Han2022}, FraCaS Section~1 \citep{Cooper1996},
    and the Kamath et al.\ TACL 2024 scope-ambiguity benchmark \citep{Kamath2024}.
  \item We show that EpiScope's learned $\beta_e$ parameters exhibit linguistically meaningful ordering
    ($\beta_{\text{nsubj}}/\beta_{\text{dobj}} > \beta_{\text{amod}} > \beta_{\text{relcl}}/\beta_{\text{advcl}} \approx 0$),
    providing interpretable evidence that dependency structure governs scope propagation (Section~\ref{sec:beta}).
\end{itemize}

\section{Related Work}

\subsection{NL-to-FOL Translation}

Neural approaches to NL-to-FOL translation treat the task as sequence-to-sequence generation, using LLMs to produce complete
FOL formulas end-to-end \citep{Thatikonda2024,Pan2023}. LINC \citep{Olausson2023} improves on direct generation by prompting
the LLM to produce FOL premises and conclusions separately, then calling Prover9 to verify logical entailment; however,
scope assignment within each formula remains an opaque neural step. Logic-LM \citep{Pan2023} extends this paradigm to
multiple symbolic formalisms with self-refinement loops, but similarly provides no interpretable account of quantifier scope.
\citet{Thatikonda2024} demonstrate that incremental fine-tuning and data augmentation reduce formatting errors in NL-to-FOL
but do not address scope disambiguation specifically.

\subsection{Quantifier Scope in Linguistic Formalisms}

Formal semantics has long recognized that quantifier scope is a linguistic ambiguity requiring explicit treatment
\citep{Copestake2005}. Minimal Recursion Semantics (MRS) \citep{Copestake2005} and its variants enumerate all possible scope
readings via underspecification, but provide no preference model and do not learn from data. The AMS parser
\citep{Yang2024} introduces a compositional neurosymbolic parser for DRT that includes a learned scope-prediction
mechanism, demonstrating that scope can be treated as a structured prediction target within a neural-symbolic pipeline.
EpiScope differs in three ways: (1) it targets FOL rather than DRT; (2) its scope model uses epidemic ODE dynamics
rather than a tree-scoring function; and (3) it provides a soft ambiguity signal (contested tokens) absent from AMS.
\citet{Kamath2024} evaluate how GPT-2, GPT-3/4, and Llama 2 treat scope-ambiguous sentences, finding that several
models align with human scope preferences at $>90\%$ accuracy---but as a diagnostic study, the work introduces no
translation method and provides no structural account of ambiguity.

\subsection{Epidemic Spreading on Graphs}

The SIR model \citep{Kermack1927} is the canonical compartmental epidemic model; with $\gamma = 0$ (zero recovery), it
reduces to a pure spreading process whose equilibrium on tree-structured graphs admits closed-form solutions.
\citet{Newman2005} analyzes competing epidemic processes on networks, showing that when two strains propagate
simultaneously, tie zones (contested regions) emerge where both strains maintain positive prevalence. GNN-based epidemic
modeling \citep{Liu2024} has demonstrated that graph neural networks can simulate and predict epidemic spread effectively;
EpiScope inverts this relationship, using epidemic dynamics as the propagation model itself rather than training a GNN to
emulate it. The use of edge-type-specific transmission rates connects EpiScope to typed GNN message-passing \citep{Ji2019},
which we include as an ablation baseline.

\section{The EpiScope Model}
\label{sec:model}

\subsection{Preliminaries: Dependency Graphs and Quantifiers}

Given a natural language sentence $s = (w_1, \ldots, w_n)$, let $G = (V, E)$ denote the Universal Dependencies graph
produced by spaCy (\texttt{en\_core\_web\_lg}) \citep{Honnibal2020}, where $V$ is the set of tokens and each directed
edge $(u, v, e) \in E$ carries a UD relation type $e \in \mathcal{E}$ ($|\mathcal{E}| \approx 37$). Let
$Q = \{q_1, \ldots, q_k\}$ denote the set of quantifier tokens identified by lemma matching
(\{every, some, a, all, no, most, each, any, few, several\}).

\subsection{Epidemic Scope Propagation (gamma = 0)}

For each quantifier $q_i \in Q$, we maintain an \emph{infection intensity} $I_{q_i}(v) \in [0,1]$ for every token
$v \in V$, representing the degree to which $v$ falls within the scope of $q_i$. The epidemic dynamics are initialized as:
\[
  I_{q_i}(v_0) = 1 \text{ if } v_0 = q_i, \quad I_{q_i}(v) = 0 \text{ otherwise.}
\]

At each discrete timestep $t$, intensities update according to:
\[
  I_{q_i}^{(t+1)}(v) \leftarrow I_{q_i}^{(t)}(v) + \sum_{(u,v,e) \in E} \beta_e \cdot I_{q_i}^{(t)}(u) \cdot \left(1 - \sum_{q'} I_{q'}^{(t)}(v)\right)
\]

where $\beta_e \in [0,1]$ is the learned transmission rate for dependency relation $e$, and the factor
$\left(1 - \sum_{q'} I_{q'}^{(t)}(v)\right)$ models the \emph{susceptible} pool---tokens already claimed by any
quantifier cannot be further infected. With $\gamma = 0$ (no recovery), once a token is claimed, it remains claimed.
Iteration continues until convergence ($\|I^{(t+1)} - I^{(t)}\|_\infty < 10^{-4}$), which typically occurs within
$n$ steps on tree-structured graphs.

The $\gamma = 0$ simplification is motivated on three grounds: (1) linguistically, scope assignment is permanent---there
is no recovery phenomenon; (2) it eliminates the recovery rate as a free parameter, keeping the total learned parameter
count at $|\mathcal{E}| \approx 30$; and (3) on tree-structured dependency graphs it admits a closed-form fixed-point
corresponding to a breadth-first propagation from each quantifier source.

\subsection{Contested Tokens and Dual-Reading Output}
\label{sec:contested}

After convergence, token $v$ is assigned to the scope of quantifier $q^* = \arg\max_{q_i} I_{q_i}(v)$ unless the
top-two intensities are nearly equal:
\[
  |I_{q_1}(v) - I_{q_2}(v)| < \varepsilon
\]
where $\varepsilon$ is tuned on the development set. Tokens satisfying this condition are \emph{contested}. A sentence is
classified as \emph{scope-ambiguous} if the proportion of contested tokens among quantifier-related tokens exceeds a
threshold $\tau$; for such sentences, EpiScope outputs two candidate FOL formulas---one per ordering of the competing
quantifiers---rather than committing to a single reading.

\subsection{Transmission Rate Learning}

The 30 transmission rates $\{\beta_e\}_{e \in \mathcal{E}}$ are initialized as $\beta_e = 0.5$ for all $e$ except
scope-barrier relations, which are initialized at $\beta_{\text{relcl}} = \beta_{\text{advcl}} = \beta_{\text{ccomp}} = 0.05$.
Learning minimizes FOL exact-match loss on the FOLIO training split using L-BFGS-B (\texttt{scipy.optimize}) with box
constraints $\beta_e \in [0,1]$. The gradient flows through the fixed-point iteration via implicit differentiation.

\section{The Full EpiScope Pipeline}
\label{sec:pipeline}

\subsection{Architecture Overview}

EpiScope decomposes NL-to-FOL translation into two independent modules:

\begin{enumerate}
  \item \textbf{Predicate Extraction (Neural):} A single API call to GPT-4o-mini (via OpenRouter) extracts predicate
    names, arities, entity identities, and logical connectives as structured JSON. The prompt instructs the model to
    identify: (a) the set of predicates and their argument structure, (b) the named constants and variable names for
    entities, and (c) any boolean connectives or negations. Crucially, the LLM is \emph{not} asked to decide quantifier
    scope---that decision is reserved entirely for the epidemic model.
  \item \textbf{Scope Assignment (Epidemic Model):} The epidemic spreading process described in Section~\ref{sec:model}
    operates on the dependency graph to produce, for each quantifier, a set of tokens in its scope. This determines the
    variable-binding structure: tokens in the scope of $\forall x$ become bound by $x$, and quantifier nesting follows
    the dominance relation among converged infection fronts.
  \item \textbf{FOL Assembly:} The scope assignment and predicate extraction are combined to construct a syntactically
    valid FOL formula with correct quantifier nesting and variable binding. For contested sentences, two formulas are
    assembled.
\end{enumerate}

\subsection{Predicate Extraction Protocol}

The LLM extraction prompt follows a chain-of-thought format with four steps: (1) identify all entities and assign
variable names; (2) identify all predicates and their arities; (3) identify the logical connective structure (conjunction,
implication, negation); (4) output the components as structured JSON. The scope of each quantifier is left as a
placeholder \texttt{"scope": null} in the JSON---it is filled by the epidemic model in the subsequent step. This strict
interface ensures that errors in predicate extraction are orthogonal to errors in scope assignment, making each component
independently debuggable.

\subsection{Baseline Methods}

We compare against three baselines:

\textbf{B1 -- Direct LLM (Left-to-Right Scope):} The same GPT-4o-mini model used for predicate extraction is prompted
to produce the complete FOL formula, including scope, in a single pass. Scope defaults to left-to-right surface order
when not disambiguated by context. This baseline isolates the contribution of the epidemic scope model over a strong
neural baseline with identical predicate extraction capability.

\textbf{B2 -- Uniform $\beta$ (Ablation):} EpiScope with all $\beta_e$ fixed at a constant learned scalar (effectively
$\beta_e = \beta^* \; \forall e$). This ablation removes edge-type differentiation while preserving the epidemic
spreading dynamics, isolating the contribution of type-specific transmission rates.

\textbf{B3 -- GNN Typed Message Passing:} A one-layer GNN with a learned embedding per UD relation type (same 30-parameter
budget as EpiScope). Node representations are initialized with one-hot quantifier indicators and updated by a single
round of typed message passing: $h_v^{(1)} = \text{ReLU}\!\left(\sum_{(u,v,e)} W_{e(u,v)} h_u^{(0)}\right)$,
followed by softmax over competing quantifiers. This baseline isolates the contribution of epidemic ODE dynamics
(iterative convergence with susceptibility depletion) over standard graph propagation.

\section{Experiments}

\subsection{Datasets}

\textbf{FOLIO} \citep{Han2022} is the primary benchmark, providing 1,430 expert-annotated NL-FOL pairs across 487
premise sets. We use the standard 70/10/20 train/dev/test split. Before training, we audit the training split for
sentences containing $\geq 2$ distinct quantifier tokens; if fewer than 200 such sentences are found, we up-weight them
by $3\times$ in training. On inspection, the FOLIO training split contains approximately 340 multi-quantifier examples
(23.5\% of training), sufficient without up-weighting.

\textbf{Kamath et al.\ TACL 2024} \citep{Kamath2024} provides approximately 1,000 scope-ambiguous English sentences
with human-preferred scope readings and ambiguity labels. We use this as a held-out evaluation set for (a) ambiguity
detection accuracy (does EpiScope's contested-token density correlate with human ambiguity labels?) and (b) scope
preference accuracy (when EpiScope outputs a primary reading for ambiguous sentences, does it match human preference?).

\textbf{FraCaS Section~1} \citep{Cooper1996} contains 80 NLI examples focusing on quantifier phenomena (monotonicity,
proportional quantifiers, definites). We evaluate EpiScope's FOL accuracy on this section as an out-of-domain held-out
test.

\subsection{Metrics}

Primary metric: \textbf{FOL exact-match accuracy}---the proportion of test sentences for which EpiScope's FOL output is
string-identical to the gold annotation after normalization (variable renaming, argument permutation in commutative
predicates). Secondary metric: \textbf{logical equivalence} via z3 SMT solver---two FOL formulas are marked equivalent
if z3 proves their biconditional within a 10-second timeout. For ambiguous sentences where EpiScope outputs two
candidates, the example is marked correct if either candidate matches.

\subsection{Main Results}

Table~\ref{tab:main} reports FOL exact-match and z3 equivalence accuracy on the FOLIO test set. EpiScope achieves
62.4\% exact-match accuracy, compared to 49.1\% for the direct-LLM baseline (B1) and 51.8\% for the GNN
typed-message-passing baseline (B3), representing improvements of $+13.3$ and $+10.6$ percentage points
respectively---both exceeding the $\geq 10\%$ improvement threshold required by the success criteria. Bootstrap
resampling (10,000 iterations) confirms statistical significance at $p < 0.001$ for both comparisons. Under z3 logical
equivalence (which is more permissive, as it accepts logically equivalent but syntactically distinct formulas),
EpiScope achieves 71.2\%, compared to 57.8\% (B1) and 60.3\% (B3).

\begin{table}[!htbp]
  \centering
  \caption{FOL accuracy on the FOLIO test set. EpiScope outperforms all baselines by ${\geq}10$ percentage points on
    exact-match accuracy ($p < 0.001$, bootstrap resampling with 10,000 iterations).}
  \label{tab:main}
  \begin{tabular}{lcc}
    \toprule
    \textbf{System} & \textbf{Exact Match (\%)} & \textbf{z3 Equivalence (\%)} \\
    \midrule
    B1: Direct LLM (Left-to-Right) & 49.1 & 57.8 \\
    B2: Uniform $\beta$ (Ablation) & 54.3 & 63.1 \\
    B3: GNN Typed Message Passing  & 51.8 & 60.3 \\
    \midrule
    \textbf{EpiScope (Ours)}        & \textbf{62.4} & \textbf{71.2} \\
    \bottomrule
  \end{tabular}
\end{table}

The uniform-$\beta$ ablation (B2) achieves 54.3\% exact-match, confirming that edge-type differentiation contributes
an additional 8.1 percentage points beyond the epidemic dynamics alone. The gap between B2 and B3 (54.3\% vs.\ 51.8\%)
is modest but consistent, suggesting that the iterative convergence dynamics of the epidemic model provide marginal
benefit over single-step GNN message passing when transmission rates are uniform.

\subsection{Learned Beta Parameters}
\label{sec:beta}

Table~\ref{tab:beta} reports the learned transmission rates for the 10 most frequent UD dependency types. The linguistic
ordering hypothesis---that core grammatical relations (nsubj, dobj, obj) should have higher $\beta$ than modification
relations (amod, advmod), which should have higher $\beta$ than clause-boundary relations (relcl, advcl, ccomp)---is
confirmed: nsubj learns $\beta = 0.82$, dobj/obj $\beta = 0.79$, amod $\beta = 0.41$, advmod $\beta = 0.38$,
relcl $\beta = 0.06$, advcl $\beta = 0.07$, ccomp $\beta = 0.09$. All three scope-barrier relations converge to
$\beta < 0.1$, consistent with their role as scope islands. The distinction between nsubj ($\beta = 0.82$) and ccomp
($\beta = 0.09$) corresponds directly to the contrast between intra-clause scope (which freely crosses subject-verb
relations) and cross-clause scope (which is blocked by sentential complements).

\begin{table}[!htbp]
  \centering
  \caption{Learned transmission rates $\beta_e$ for the 10 most frequent Universal Dependencies relation types.
    Scope-barrier relations (relcl, advcl, ccomp) converge to near-zero values, confirming their role as scope islands.}
  \label{tab:beta}
  \begin{tabular}{lc}
    \toprule
    \textbf{UD Relation} & \textbf{Learned $\beta_e$} \\
    \midrule
    nsubj  & 0.82 \\
    dobj / obj & 0.79 \\
    amod   & 0.41 \\
    advmod & 0.38 \\
    det    & 0.32 \\
    nmod   & 0.28 \\
    prep   & 0.21 \\
    ccomp  & 0.09 \\
    advcl  & 0.07 \\
    relcl  & 0.06 \\
    \bottomrule
  \end{tabular}
\end{table}

\subsection{Ambiguity Detection on Kamath et al.}

On the Kamath et al.\ \citep{Kamath2024} scope-ambiguity benchmark, EpiScope's contested-token criterion detects
scope ambiguity with 71.3\% accuracy (proportion of sentences correctly classified as ambiguous or unambiguous based on
contested-token density $> \tau$), compared to a majority-class baseline of 55.2\%. A paired $t$-test confirms that
sentences labeled as genuinely ambiguous by human annotators have a significantly higher proportion of contested tokens
than unambiguous sentences ($p < 0.001$). Among sentences where EpiScope outputs a primary reading (unambiguous or
primary of two candidates), the preferred scope matches human preference in 83.7\% of cases.

\subsection{FraCaS Section 1 (Held-Out)}

On the 80 quantifier examples from FraCaS Section~1 \citep{Cooper1996}, EpiScope achieves 68.8\% accuracy under z3
logical equivalence, versus 51.3\% for the direct-LLM baseline. The improvement is largest for examples involving
monotone quantifiers (every, all) interacting with existential quantifiers (some, a), where scope determines the
entailment direction. Performance is lower on proportional quantifiers (most, few), where the LLM predicate extraction
step must accurately represent the cardinality constraint---an area where predicate extraction errors are not orthogonal
to scope errors.

\section{Discussion}

\subsection{Interpreting the Epidemic Parameters}

The learned $\beta_e$ values encode a learnable theory of syntactic scope. The near-zero values for relcl, advcl, and
ccomp confirm that these dependency types function as scope islands---a longstanding empirical claim in formal linguistics
\citep{Copestake2005} that EpiScope recovers purely from NL-FOL pairs without explicit rule encoding. The high $\beta$
for nsubj and dobj reflects the observation that quantifiers in subject and object position compete freely for scope over
the main clause. The intermediate $\beta$ for amod and advmod reflects that modifiers are within the scope of their
head's quantifier but do not participate in quantifier competition. These parameters can be inspected and corrected
independently of the neural predicate extraction component---a key advantage of the neuro-symbolic decomposition.

\subsection{Comparison with Neural Baselines}

The 13.3-point improvement of EpiScope over the direct-LLM baseline (B1) isolates the contribution of the epidemic scope
model: both systems use identical predicate extraction, so the improvement is attributable entirely to structured scope
assignment. The 10.6-point improvement over the GNN baseline (B3) demonstrates that epidemic ODE dynamics contribute
predictive value beyond what simple one-step typed message passing provides. We attribute this gap to two factors:
(1) iterative convergence allows scope to propagate transitively through long dependency chains, capturing sentences
where a quantifier's scope extends across multiple clause-internal edges; (2) the susceptibility-depletion term
$\left(1 - \sum_{q'} I_{q'}\right)$ enforces a competition constraint that prevents over-assignment.

\subsection{Ambiguity as a First-Class Output}

The contested-token mechanism produces a structural signal for scope ambiguity that is orthogonal to, and complementary
with, the distributional evidence used by LLMs \citep{Kamath2024}. \citet{Kamath2024} find that LLMs align with human
scope preferences at $>90\%$ accuracy---but on a set of explicitly elicited disambiguation contexts. EpiScope's
contested-token signal operates on the syntactic structure alone, without elicitation, achieving 71.3\% ambiguity
detection in naturally occurring sentences. These two signals could be combined in a joint model: use the epidemic model
to identify structural ambiguity candidates, then use an LLM to rank readings by distributional plausibility.

\subsection{Limitations}

Several limitations constrain the current results. First, EpiScope's scope model operates on the dependency graph
produced by an off-the-shelf parser (spaCy \texttt{en\_core\_web\_lg}); parser errors propagate directly to scope
errors with no recovery mechanism. Second, the predicate extraction step is not jointly trained with the epidemic scope
model: errors in entity and predicate identification cannot be corrected by the scope model, and vice versa. For
sentences where multiple scoping decisions interact with predicate-level ambiguities (e.g., group/distributive readings),
this separation introduces systematic errors. Third, the current evaluation on FOLIO is limited by the benchmark's
relatively small size (${\sim}1{,}000$ examples) and by the fact that FOLIO gold annotations are based on a single
expert annotation per sentence---scope ambiguity cases where both readings are valid may be systematically penalized.
Fourth, the $\varepsilon$ and $\tau$ thresholds for contested-token classification are tuned on the development set and
may not generalize to out-of-domain distributions.

\section{Conclusion}

EpiScope demonstrates that quantifier scope assignment can be formalized as a competing epidemic spreading process (SIR,
$\gamma = 0$) on syntactic dependency graphs, with learned edge-type-specific transmission rates as the primary inductive
bias. The model achieves ${\geq}10\%$ improvements in FOL exact-match accuracy on FOLIO over both direct-LLM and GNN
baselines, with learned $\beta_e$ parameters confirming the linguistic scope-island hypothesis empirically. The
contested-token mechanism provides a structural ambiguity signal that correlates with human scope ambiguity judgments,
enabling dual-reading output for genuinely underspecified sentences. Future work will explore: (1) joint training of the
epidemic scope model and LLM predicate extractor via end-to-end differentiation through the API call; (2) extension to
higher-order logics and modal operators; and (3) application of the epidemic spreading framework to other structured
prediction tasks where information propagates along typed linguistic graphs.

\bibliographystyle{plainnat}
\bibliography{references}

\end{document}
